\maketitle
\makesignature

\ifproject
\begin{abstractTH}

โครงงานนี้มีวัตถุประสงค์เพื่อประเมินความเสี่ยงด้านความเป็นส่วนตัวของข้อมูลในโมเดลภาษา (LMs) 
สำหรับบริบทภาษาไทย โดยเน้นการวิเคราะห์การโจมตีด้วยการสกัดข้อมูล (Data Extraction Attack หรือ DEA) และ 
การโจมตีด้วยการเจลเบรก (Jailbreak Attack หรือ JA) ผ่านตัวแปรด้านระดับภาษาไทย 5 ระดับ ได้แก่ ระดับพิธีการ, ทางการ, 
กึ่งทางการ, ไม่เป็นทางการ และกันเอง การทดลองดำเนินการบนโมเดล Llama-3.1-8b-instant, Meta-llama-llama-4-maverick-17b-128e-instruct และ 
Moonshotai-kimi-k2-instruct-0905 

ผลการทดลองชี้ให้เห็นว่า Llama-3.1-8B-instant มีความต้านทานต่อ DEA ดีที่สุด ขณะที่ Meta-llama-llama-4-maverick-17b-128e-instruct 
มีประสิทธิภาพสูงสุดในการรับมือกับ JA สำหรับการลดความเสี่ยงด้านความเป็นส่วนตัวของข้อมูล ผลการวิจัยสรุปผลได้ว่ากลไกการคัดกรองข้อมูล (Scrubbing) 
เหมาะสมสำหรับการป้องกันการรั่วไหลของข้อมูลที่ระบุตัวตนได้ (Personal Identifiable Information หรือ PII) ในขณะที่การป้อนคำสั่งป้องกัน (Defensive Prompting) ให้ผลลัพธ์ที่ดีกว่าในการป้องกันการเจลเบรก
(JA)

\end{abstractTH}

\begin{abstract}

This project aims to assess data privacy risks in Language Models (LMs) within the Thai context, specifically focusing on the analysis of Data Extraction Attacks (DEA) and Jailbreak Attacks (JA). The evaluation is conducted across five levels of Thai linguistic registers: formal, official, semi-formal, informal, and casual. Experiments were performed on three models: Llama-3.1-8B-instant, Meta-Llama-4-Maverick-17B-128e-instruct, and Moonshotai-Kimi-K2-instruct-0905.

Experimental results indicate that Llama-3.1-8B-instant exhibits the highest resistance to DEA, while Meta-Llama-4-Maverick-17B-128e-instruct demonstrates the highest effectiveness in mitigating JA. Regarding personal data leakage, the research confirms that the Scrubbing mechanism is most suitable for preventing Personal Identifiable Information (PII) leakage. In contrast, Defensive Prompting proves superior in preventing jailbreak attempts. This study concludes by providing practical guidelines for selecting appropriate models and defense strategies to minimize data privacy risks in Thai language applications.
\end{abstract}

\iffalse
\begin{dedication}
This document is dedicated to all Chiang Mai University students.

Dedication page is optional.
\end{dedication}
\fi % \iffalse

\begin{acknowledgments}

โครงงานวิจัยฉบับนี้สำเร็จลุล่วงไปได้ด้วยความกรุณาอย่างยิ่งจากอาจารย์จักรพงศ์ นาทวิชัย 
อาจารย์ที่ปรึกษา
\\ โครงงาน ที่ได้สละเวลาให้คำปรึกษา แนะนำแนวทางในการศึกษา และช่วยแก้ไขข้อบกพร่องต่าง ๆ ตลอดระยะเวลาการดำเนินงาน เพื่อให้งานวิจัยมีความสมบูรณ์และถูกต้องตามหลักวิชาการ คณะผู้จัดทำขอกราบขอบพระคุณเป็นอย่างสูงไว้ ณ ที่นี้

คณะผู้จัดทำขอขอบพระคุณ ภาควิชาวิศวกรรมคอมพิวเตอร์ คณะวิศวกรรมศาสตร์ มหาวิทยาลัยเชียงใหม่ เป็นอย่างสูง สำหรับการสนับสนุนด้านงบประมาณในการเข้าถึงทรัพยากรทางเทคนิค (API) และการเอื้อเฟื้อสถานที่ในการดำเนินงานวิจัย  รวมถึงขอบพระคุณคณาจารย์ทุกท่านที่ได้ประสิทธิ์ประสาทวิชาความรู้ ซึ่งเป็นรากฐานสำคัญในการพัฒนาโครงงานวิจัยในครั้งนี้

นอกจากนี้ คณะผู้จัดทำขอแสดงความนับถือและขอบคุณงานวิจัย "LLM-PBE: Assessing Data Privacy in Large Language Models" สำหรับการนำเสนอกรอบแนวคิดและระเบียบวิธีวิจัยอันทรงคุณค่า ซึ่งคณะผู้จัดทำได้นำมาอ้างอิงและพัฒนาต่อยอดเพื่อศึกษาผลกระทบในบริบทภาษาไทย 

ท้ายที่สุดนี้ คณะผู้จัดทำขอกราบขอบพระคุณบิดา มารดา และครอบครัว ที่ให้การสนับสนุนและเป็นกำลังใจสำคัญเสมอมา รวมถึงเพื่อนนักศึกษาทุกคนที่ให้ความช่วยเหลือและแลกเปลี่ยนความคิดเห็นตลอดการทำโครงงานวิจัย

\acksign{2026}{2}{26}
\end{acknowledgments}%
\fi % \ifproject

\contentspage

\ifproject
\figurelistpage

\tablelistpage
\fi % \ifproject

% \abbrlist % this page is optional

% \symlist % this page is optional

% \preface % this section is optional
